\subsection*{序章で使う案内語}

この一覧は、序章の地図を読むための案内です。
技術的な仕組みは担当章で詳しく説明します。

\begin{description}
  \item[\term{LLM}] Large Language Modelの略で、大量の文から言語の規則性を学ぶモデルです。
  \item[\term{計算}] 入力を定めた規則で別の値へ変換する処理です。
  \item[\term{研修}] 知識と技能を、説明、演習、確認によって身に付ける学習活動です。
  \item[\term{読者}] 本書を読み、演習を行う新卒アプリケーションエンジニアです。
  \item[\term{実務}] 実際の開発、運用、調査、説明で行う仕事です。
  \item[\term{現在地}] 一要求が、九つの場面のどこにあるかを示す案内です。
  \item[\term{旅}] 一要求の表現が層を移る流れを、本書で説明するための比喩です。
  \item[\term{注文確認AI}] 注文情報を取得し、利用者向けの文を生成する本書の例題アプリです。
  \item[\term{図1}] 序章の図\ref{fig:ch00-overview}を指すカタログ上の旧称です。
  \item[\term{前方参照}] これから説明する章や節を先に案内する参照です。
  \item[\term{後方参照}] すでに説明した章や節へ戻る参照です。
  \item[\term{演習}] 読者が操作、計算、観測を行う課題です。
  \item[\term{確認問題}] 用語と章の意図を説明できるか確かめる問いです。
\end{description}
